% ======================================================================
% Appendices A--D (proofs).  Uses the \providecommand shorthands of sections_II_IV.tex.
\appendix
\section*{Appendix: Complete Mathematical Proofs}

\subsection*{A. Proof of Proposition~\ref{prop:twin_crown} (Coupled Twin-Network CROWN)}
Let \(y_\ell(z)\) and \(y_\ell^q(z)\) denote the pre-activation vectors at layer \(\ell\) of \(\mathcal{N}\) and \(\mathcal{N}_q\) evaluated at the port input \(v = g_r^TPz\), with post-activations \(s_\ell := \sigma_{\mathrm{GS2}}(y_\ell)\), \(s_\ell^q := \sigma_{\mathrm{GS2}}(y_\ell^q)\) (\(s_0 := v\), \(s_0^q := v + e_0\)) and difference \(\delta y_\ell(z) := y_\ell^q(z) - y_\ell(z)\). Write \(\Delta W_\ell := W_\ell^q - W_\ell\) and \(\Delta b_\ell := b_\ell^q\) (the teacher biases vanish). At \(\ell = 1\),
\(\delta y_1(z) = \Delta W_1 v + W_1^q e_0 + \Delta b_1 + e_1\) is an exact affine function of \((z, e_0, e_1)\), because \(v = g_r^TPz\) is linear in \(z\). By the shared ReLU routing identity \eqref{eq:groupsort_relu}, the post-activation difference for each pair \((2i-1, 2i)\) is
\begin{equation}
\label{eq:gs2_diff_pair}
\sigma_{\mathrm{GS2}}(y_{\ell,2i-1}^q, y_{\ell,2i}^q) - \sigma_{\mathrm{GS2}}(y_{\ell,2i-1}, y_{\ell,2i}) = \begin{bmatrix} \delta y_{\ell,2i}(z) + \nu_{\ell,i}(z) \\ \delta y_{\ell,2i-1}(z) - \nu_{\ell,i}(z) \end{bmatrix},
\end{equation}
where \(\nu_{\ell,i}(z) := \mathrm{ReLU}(d_{\ell,i}(z) + \delta d_{\ell,i}(z)) - \mathrm{ReLU}(d_{\ell,i}(z))\) and \(\delta d_{\ell,i}(z) = \delta y_{\ell,2i-1}(z) - \delta y_{\ell,2i}(z) \in [\underline{\delta}_{\ell,i}, \overline{\delta}_{\ell,i}]\). On any sub-hyperrectangle \(\mathcal{Z}_s \subset \ZO\):
\begin{enumerate}
    \item For \(i \in \mathcal{I}_\ell^+(\mathcal{Z}_s)\) (\(\underline{d}_{\ell,i} \ge 0\) and \(\underline{d}_{\ell,i}^q \ge 0\)), \(\nu_{\ell,i}(z) = 1 \cdot \delta d_{\ell,i}(z) + 0\) identically on \(\mathcal{Z}_s\) (\(\alpha_{\ell,i}^L = \alpha_{\ell,i}^U = 1\), \(\beta_{\ell,i}^L = \beta_{\ell,i}^U = 0\)).
    \item For \(i \in \mathcal{I}_\ell^-(\mathcal{Z}_s)\) (\(\overline{d}_{\ell,i} \le 0\) and \(\overline{d}_{\ell,i}^q \le 0\)), \(\nu_{\ell,i}(z) = 0 \cdot \delta d_{\ell,i}(z) + 0\) identically on \(\mathcal{Z}_s\) (\(\alpha_{\ell,i}^L = \alpha_{\ell,i}^U = 0\), \(\beta_{\ell,i}^L = \beta_{\ell,i}^U = 0\)).
    \item For \(i \in \mathcal{I}_\ell^*(\mathcal{Z}_s)\) (sign-indeterminate, including any comparator where quantization flips the sorting order between \(d_{\ell,i}\) and \(d_{\ell,i}^q\)), monotonicity and \(1\)-Lipschitz continuity of \(\mathrm{ReLU}\) imply the universal pointwise envelope
    \[
    \min(0, \delta d_{\ell,i}) = -\mathrm{ReLU}(-\delta d_{\ell,i}) \le \nu_{\ell,i}(z) \le \mathrm{ReLU}(\delta d_{\ell,i}) = \max(0, \delta d_{\ell,i}) \quad \forall d_{\ell,i} \in \mathbb{R},
    \]
    whose upper envelope is attained at \(d_{\ell,i}(z) = 0\) (where \(\nu_{\ell,i} = \mathrm{ReLU}(\delta d_{\ell,i})\)) and whose lower envelope is attained at \(d_{\ell,i}(z) = -\delta d_{\ell,i}\), so both sides are tight. If \(\underline{\delta}_{\ell,i} \ge 0\), then \(\min(0, \delta d_{\ell,i}) = 0\) and \(\max(0, \delta d_{\ell,i}) = \delta d_{\ell,i}\) are already exact linear functions; similarly, if \(\overline{\delta}_{\ell,i} \le 0\), then \(\min(0, \delta d_{\ell,i}) = \delta d_{\ell,i}\) and \(\max(0, \delta d_{\ell,i}) = 0\) are exact linear functions. When \(\underline{\delta}_{\ell,i} < 0 < \overline{\delta}_{\ell,i}\), using \(\max(0, \delta d_{\ell,i}) = \mathrm{ReLU}(\delta d_{\ell,i})\) and \(\min(0, \delta d_{\ell,i}) = \delta d_{\ell,i} - \mathrm{ReLU}(\delta d_{\ell,i})\) and bounding \(\mathrm{ReLU}(\delta d_{\ell,i})\) from above by its secant chord \(\alpha_{\ell,i}^U \delta d_{\ell,i} + \beta_{\ell,i}^U\) (exact at the endpoints \(\underline{\delta}_{\ell,i}, \overline{\delta}_{\ell,i}\), and an upper bound in between by convexity of \(\mathrm{ReLU}\)) and from below by \(0\) yields both the affine upper bound \(\alpha_{\ell,i}^U \delta d_{\ell,i} + \beta_{\ell,i}^U\) and the affine lower bound \(\delta d_{\ell,i} - (\alpha_{\ell,i}^U \delta d_{\ell,i} + \beta_{\ell,i}^U) = (1 - \alpha_{\ell,i}^U)\delta d_{\ell,i} - \beta_{\ell,i}^U = \alpha_{\ell,i}^L \delta d_{\ell,i} + \beta_{\ell,i}^L\), establishing \eqref{eq:twin_affine_relax}--\eqref{eq:twin_slopes}.
\end{enumerate}

\textbf{Layer propagation and part (ii).} At layer \(\ell+1\), decompose
\(\delta y_{\ell+1}(z) = \Delta W_{\ell+1}\sigma_{\mathrm{GS2}}(y_\ell(z)) + \Delta b_{\ell+1} + e_{\ell+1} + W_{\ell+1}^q [\sigma_{\mathrm{GS2}}(y_\ell^q(z)) - \sigma_{\mathrm{GS2}}(y_\ell(z))]\). The first term is the single teacher network multiplied by the fixed matrix \(\Delta W_{\ell+1}\); it is relaxed by single-network CROWN secant bounds on \([\underline{d}_{\ell,i}, \overline{d}_{\ell,i}]\) for teacher comparators that are not sign-stable on \(\mathcal{Z}_s\), and is exactly affine otherwise. The terms \(\Delta b_{\ell+1}\) and \(e_{\ell+1}\) are constants and box-bounded injections, respectively. In the last term, substitute \eqref{eq:gs2_diff_pair}: by \eqref{eq:gs2_diff_pair}, the coefficient of \(\nu_{\ell,i}\) in any row of \(\Lambda W_{\ell+1}^q[\sigma_{\mathrm{GS2}}(y^q_\ell) - \sigma_{\mathrm{GS2}}(y_\ell)]\), where \(\Lambda\) is the coefficient matrix accumulated by the backward pass from the output, is \(c_{\ell,i} := [\Lambda W^q_{\ell+1}]_{\cdot,2i-1} - [\Lambda W^q_{\ell+1}]_{\cdot,2i}\). For the upper bound of a row, replace \(\nu_{\ell,i}\) by \(\alpha^U\delta d_{\ell,i} + \beta^U\) where \(c_{\ell,i} \ge 0\) and by \(\alpha^L\delta d_{\ell,i} + \beta^L\) where \(c_{\ell,i} < 0\); for the lower bound reverse the choice. By \eqref{eq:twin_affine_relax} each replacement is sound (multiplying an inequality by a non-negative, resp.\ negative, coefficient preserves, resp.\ reverses, it), and since \(\delta d_{\ell,i}\) is itself a linear combination of \(\delta y_\ell\), the substitution yields affine upper and lower bounds in \((z, e_0, \dots, e_L)\) that are recursively expressed through layer \(\ell\). Note that \(\alpha_{\ell,i}^L, \alpha_{\ell,i}^U \in [0, 1]\) and \(|\beta_{\ell,i}^{U/L}| = \frac{-\underline{\delta}_{\ell,i}\overline{\delta}_{\ell,i}}{\overline{\delta}_{\ell,i} - \underline{\delta}_{\ell,i}} \le \min(|\underline{\delta}_{\ell,i}|, \overline{\delta}_{\ell,i}) \le \overline{\delta}_{\ell,i} - \underline{\delta}_{\ell,i}\). Iterating to \(\ell = 1\), where everything is affine, and concretizing the resulting affine forms over \(\mathcal{Z}_s\) and the injection boxes (exact maximization of an affine function over a product of boxes) gives \([\Delta_{\mathrm{net}}(z)]_j \in [\underline{D}_{s,j}, \overline{D}_{s,j}]\) for all \(z \in \mathcal{Z}_s\) and all admissible injections. Hence \(\|\Delta_{\mathrm{net}}(z)\|_2^2 = \sum_j [\Delta_{\mathrm{net}}(z)]_j^2 \le \sum_j \max(|\underline{D}_{s,j}|^2, |\overline{D}_{s,j}|^2)\) on \(\mathcal{Z}_s\); since the boxes cover \(\ZO\), taking the maximum over \(s\) bounds the supremum over \(\ZO\). Finally \(\|\Psi_B(z) - \Tt(z)\|_2 = \|\Pi_U(\mathcal{N}_q(v)) - \Pi_U(\mathcal{N}(v))\|_2 \le \|\Delta_{\mathrm{net}}(z)\|_2\) by non-expansiveness of the metric projection \(\Pi_U\) onto the closed convex set \(U_{\mathrm{adm}}\). This proves \eqref{eq:twin_crown_bound}.

\textbf{Part (iii).} Let \(\delta_\ell := \sup\|\delta y_\ell\|_2\) over \(z \in \ZO\) and admissible injections, and \(\delta_0 := \sup\|s_0^q - s_0\|_2 = \sup\|e_0\|_2 \le \sqrt{H_0}\rho_0\). For \(1 \le \ell \le L\),
\(\delta y_\ell = \Delta W_\ell s_{\ell-1} + W_\ell^q(s_{\ell-1}^q - s_{\ell-1}) + \Delta b_\ell + e_\ell\). Since \(\sigma_{\mathrm{GS2}}\) is norm-preserving and \(\|W_k\|_2 \le 1\) for \(k < L\), \(\|s_{\ell-1}\|_2 \le \|v\|_2 \le \|g_r^TP\|_2\|z\|_2 \le Y\) for \(\ell - 1 \le L - 1\) (using \(\|z\|_2 \le \sqrt{K}\DO/2\) on \(\ZO\)); since \(\sigma_{\mathrm{GS2}}\) is non-expansive, \(\|s_{\ell-1}^q - s_{\ell-1}\|_2 \le \|\delta y_{\ell-1}\|_2\) for \(\ell \ge 2\). The triangle inequality and \(\|e_\ell\|_2 \le \sqrt{H_\ell}\rho_\ell\) (box \(\subset\) ball) give \(\delta_\ell \le \Delta_{W,\ell}Y + \Delta_{b,\ell} + \sqrt{H_\ell}\rho_\ell + \|W_\ell^q\|_2\,\delta_{\ell-1}\). Unrolling this linear recursion from \(\delta_0\) to \(\delta_L\) and noting \(\Delta_{\mathrm{net}} = \delta y_L\) (no output activation) yields \(\sup\|\Delta_{\mathrm{net}}\|_2 \le E_q(b_B)\) as in \eqref{eq:crown_explicit}; \(\|W_k^q\|_2 \le \|W_k\|_2 + \Delta_{W,k}\) gives the product bound. The entrywise bound \(|[\Delta W_\ell]_{ij}| \le s_\ell2^{-b_B}\) gives \(\|\Delta W_\ell\|_2 \le \|\Delta W_\ell\|_F \le \sqrt{H_\ell H_{\ell-1}}s_\ell2^{-b_B}\). The relation \(\eC \le C_{\mathrm{relax}}E_q(b_B)\) holds by definition of \(C_{\mathrm{relax}}\) (set \(C_{\mathrm{relax}} := 1\) if \(E_q = 0\), in which case all perturbations vanish and \(\eC = 0\)); its content is the bound in the exact case: if every comparator of both networks is sign-stable on \(\mathcal{Z}_s\), all relaxations are exact, \(\Delta_{\mathrm{net}}\) is exactly affine on \(\mathcal{Z}_s \times\) (injection boxes), and the concretization returns its exact component ranges, so \(\max(|\underline{D}_{s,j}|, |\overline{D}_{s,j}|) = \sup|[\Delta_{\mathrm{net}}]_j| \le \sup\|\Delta_{\mathrm{net}}\|_2 \le E_q(b_B)\), whence \(\eC \le \sqrt{m}\,E_q(b_B)\). In general \(C_{\mathrm{relax}} = 1\) is false: for \(m = 2\), a \(45^\circ\) rotation applied to an input box \([-\rho,\rho]^2\) has exact \(\ell_2\) range \(\sqrt2\rho\), while the component-wise \(\ell_2\) assembly returns \(2\rho\).

\textbf{Part (iv).} Let \(\kappa := 2^{-b_B} \le 1\). All perturbation sources are \(O(\kappa)\): \(\|\Delta W_\ell\|_2 \le \sqrt{H_\ell H_{\ell-1}}s_\ell\kappa\), \(\Delta_{b,\ell}, \rho_\ell \le c\kappa\). For \(E_q\) the claim is immediate from \eqref{eq:crown_explicit} since the products are bounded by \((L_{\mathcal{T}} + \sqrt{mH}s_L)\prod_{k<L}(1 + Hs_k)\) uniformly in \(\kappa \le 1\). For \(\eC\), argue by induction on the layer index that every interval half-width and every endpoint produced by the twin backward pass is at most \(C_\ell\kappa\) with \(C_\ell\) independent of \(\kappa\) and of the partition: the backward coefficient matrices are products of \(W_k^q\) (entries bounded uniformly in \(\kappa \le 1\)) and diagonal slope matrices with entries in \([0,1]\), hence bounded by a \(\kappa\)-independent constant; the quantities they multiply are \(\Delta W_k s_{k-1}\) (\(\le \|\Delta W_k\|_2 Y\)), \(\Delta b_k\), \(e_k\), and offsets \(\beta\) with \(|\beta_{k,i}| \le \overline{\delta}_{k,i} - \underline{\delta}_{k,i} \le 2C_k\kappa\) by the induction hypothesis (the map \((\underline\delta, \overline\delta) \mapsto (\alpha, \beta)\) is positively homogeneous of degree \(0\) in \(\alpha\) and degree \(1\) in \(\beta\)). Summing finitely many such terms gives \(C_{\ell+1}\kappa\), and the output assembly is \(\le \sqrt{m}C_L\kappa\). \hfill \(\blacksquare\)

\subsection*{B. Proof of Proposition~\ref{prop:dahlquist} (\(P\)-Weighted Matrix Measure Contraction)}
\textbf{Part 1 (\(R_0 \succ 0\)):} Let \(V(e) = \frac{1}{2}e^T P_0 e\) and define \(\varpi := P_0^{1/2} e\), so \(V(e) = \frac{1}{2}\|\varpi\|_2^2\) and \(P_0 e = P_0^{1/2} \varpi\). Differentiating along \(\dot{e} = (J_0 - R_0)P_0 e\) gives
\[
\dot{V}(e) = (P_0 e)^T J_0 (P_0 e) - (P_0 e)^T R_0 (P_0 e).
\]
Since \(J_0^T = -J_0\), the scalar \((P_0 e)^T J_0 (P_0 e)\) equals its own negative, hence \((P_0 e)^T J_0 (P_0 e) = 0\) identically for all \(e \in \mathbb{R}^K\). For the dissipative term,
\[
(P_0 e)^T R_0 (P_0 e) \ge \lambda_{\min}(R_0)\,\|P_0 e\|_2^2 = \lambda_{\min}(R_0)\,\varpi^T P_0 \varpi \ge \lambda_{\min}(R_0)\,\lambda_{\min}(P_0)\,\|\varpi\|_2^2 = 2\,\lambda_{\min}(R_0)\,\lambda_{\min}(P_0)\,V(e),
\]
proving \eqref{eq:mu_p0}.

\textbf{Part 2 (\(R_0 = \mathrm{blkdiag}(0, \gamma I_{K_m})\)):} Here \(A_{r,0} = (J_0 - R_0)P_0 = \begin{bsmallmatrix} 0 & I \\ -\Omega_r^2 & -\gamma I \end{bsmallmatrix}\). For \(P_\epsilon = \begin{bsmallmatrix} \Omega_r^2 & \epsilon I \\ \epsilon I & I \end{bsmallmatrix}\), the Schur complement \(\Omega_r^2 - \epsilon^2 I \succ 0\) holds for all \(\epsilon \in (0, \omega_{\min})\). Direct block multiplication of \(Q_\epsilon := -(P_\epsilon A_{r,0} + A_{r,0}^T P_\epsilon)\) cancels the off-diagonal \(\pm \Omega_r^2\) blocks and yields
\[
Q_\epsilon = \begin{bmatrix} 2\epsilon \Omega_r^2 & \epsilon \gamma I \\ \epsilon \gamma I & 2(\gamma - \epsilon)I \end{bmatrix},
\]
which decouples mode-by-mode into \(2 \times 2\) symmetric blocks \(Q_{\epsilon,i} = \begin{bsmallmatrix} 2\epsilon\omega_i^2 & \epsilon\gamma \\ \epsilon\gamma & 2(\gamma - \epsilon) \end{bsmallmatrix}\) for \(i = 1, \dots, K_m\). By Sylvester's criterion, \(Q_{\epsilon,i} \succ 0\) if and only if \(2\epsilon\omega_i^2 > 0\) and
\[
\det(Q_{\epsilon,i}) = \epsilon\Big(4\gamma\omega_i^2 - \epsilon\big(\gamma^2 + 4\omega_i^2\big)\Big) > 0 \iff 0 < \epsilon < \frac{4\gamma\omega_i^2}{\gamma^2 + 4\omega_i^2}.
\]
Because \(f(\omega^2) = \frac{4\gamma\omega^2}{\gamma^2 + 4\omega^2}\) is strictly increasing in \(\omega^2 > 0\), the intersection over \(i = 1, \dots, K_m\) is governed by \(\omega_{\min} = \min_i \omega_i\) and satisfies \(\frac{4\gamma\omega_{\min}^2}{\gamma^2 + 4\omega_{\min}^2} \le \frac{4\gamma\omega_{\min}^2}{4\gamma\omega_{\min}} = \omega_{\min}\), so both \(P_\epsilon \succ 0\) and \(Q_\epsilon \succ 0\) hold for all \(\epsilon \in \big(0,\, \frac{4\gamma \omega_{\min}^2}{\gamma^2 + 4\omega_{\min}^2}\big)\).

\emph{The choice \(\epsilon = \gamma/2\).} \(\frac{\gamma}{2} < \frac{4\gamma\omega_{\min}^2}{\gamma^2 + 4\omega_{\min}^2} \iff \gamma^2 + 4\omega_{\min}^2 < 8\omega_{\min}^2 \iff \gamma < 2\omega_{\min}\) (divide by \(\gamma/2 > 0\)). At \(\epsilon = \gamma/2\), \(Q_{\gamma/2} = \begin{bsmallmatrix} \gamma\Omega_r^2 & (\gamma^2/2) I \\ (\gamma^2/2) I & \gamma I \end{bsmallmatrix} = \gamma P_{\gamma/2}\) (Pencil-Rate identity). Hence for every \(e \neq 0\), \(e^T(P_{\gamma/2}A_{r,0} + A_{r,0}^TP_{\gamma/2})e/(2e^TP_{\gamma/2}e) = -\gamma/2\), so the supremum in the definition of \(\mu_{2,P_{\gamma/2}}\) equals \(-\gamma/2\), and \(\frac{d}{dt}\|e\|_{P_{\gamma/2}}^2 = -e^TQ_{\gamma/2}e = -\gamma\|e\|_{P_{\gamma/2}}^2\) integrates to \eqref{eq:mu_pe}. Since \(P_\epsilon\) and \(Q_\epsilon\) do not depend on any reduced quantity other than \(\Omega_r\), the same identity holds on the full-order momentum-damped lattice with \(\Omega_{\mathrm{full}}\) in place of \(\Omega_r\).

\emph{Metric sandwich.} For each mode, \(2|q_ip_i| \le \omega_iq_i^2 + p_i^2/\omega_i\) (AM--GM), so \(|2\epsilon q_ip_i| \le \frac{\epsilon}{\omega_i}(\omega_i^2q_i^2 + p_i^2) \le \frac{\epsilon}{\omega_{\min}}(\omega_i^2q_i^2 + p_i^2)\). Summing over \(i\), \(e^TP_\epsilon e = e^TP_0e + 2\epsilon q^Tp \in [(1 - \epsilon/\omega_{\min}), (1 + \epsilon/\omega_{\min})]\,e^TP_0e\). \hfill \(\blacksquare\)

\begin{lemma}[Pencil rate]
\label{lem:pencil_rate}
For one mode with \(\omega > 0\), \(\gamma > 0\), \(\epsilon \in (0,\omega)\), \(A_\omega = \begin{bsmallmatrix} 0 & 1 \\ -\omega^2 & -\gamma \end{bsmallmatrix}\), \(P_{\epsilon,\omega} = \begin{bsmallmatrix} \omega^2 & \epsilon \\ \epsilon & 1\end{bsmallmatrix}\) and \(Q = -(P_{\epsilon,\omega}A_\omega + A_\omega^TP_{\epsilon,\omega})\), the largest \(\lambda\) with \(Q \succeq 2\lambda P_{\epsilon,\omega}\) is
\(r(\epsilon,\omega) = \frac{\gamma}{2} - \frac{\omega|\gamma - 2\epsilon|}{2\sqrt{\omega^2 - \epsilon^2}} \le \frac{\gamma}{2}\), with equality iff \(\epsilon = \gamma/2\).
\end{lemma}
\begin{proof}
With \(t := \gamma/2 - \lambda\) and \(a := \gamma/2 - \epsilon\), \(Q - 2\lambda P_{\epsilon,\omega} = \begin{bsmallmatrix} 2\omega^2(t - a) & 2\epsilon t \\ 2\epsilon t & 2(t + a)\end{bsmallmatrix}\) (entry by entry: \(2\epsilon\omega^2 - 2\lambda\omega^2\), \(\epsilon\gamma - 2\lambda\epsilon\), \(2(\gamma - \epsilon) - 2\lambda\)), with determinant \(4[\omega^2(t^2 - a^2) - \epsilon^2t^2] = 4[(\omega^2 - \epsilon^2)t^2 - \omega^2a^2]\). A symmetric \(2\times2\) matrix is PSD iff its diagonal entries and determinant are non-negative, i.e.\ iff \(t \ge |a|\) and \(t \ge \omega|a|/\sqrt{\omega^2 - \epsilon^2}\); the second implies the first since \(\omega/\sqrt{\omega^2 - \epsilon^2} \ge 1\). Hence \(\lambda \le \gamma/2 - \omega|a|/\sqrt{\omega^2-\epsilon^2} = r(\epsilon,\omega)\).
\end{proof}

\begin{lemma}[Characterization of (H2a) in class (H)]
\label{lem:h2a_char}
Let \(\Phi = \mathrm{blkdiag}(\Phi_q, \Phi_q)\), \(\Phi_q^T\Phi_q = I\), \(S := \Phi_q^T\Omega_{\mathrm{full}}^2\Phi_q\). Then \(P_\perp = \mathrm{blkdiag}(I - \Phi_qS^{-1}\Phi_q^T\Omega_{\mathrm{full}}^2,\ I - \Phi_q\Phi_q^T)\), and \(\Phi^T\Pfe P_\perp\) has zero diagonal blocks, zero \((1,2)\) block, and \((2,1)\) block \(\epsilon(\Phi_q^T - S^{-1}\Phi_q^T\Omega_{\mathrm{full}}^2)\). Hence (H2a) holds iff \(\Omega_{\mathrm{full}}^2\Phi_q = \Phi_qS\), i.e.\ iff \(\mathrm{Im}(\Phi_q)\) is \(\Omega_{\mathrm{full}}^2\)-invariant; in that case \(\Pf\Phi = \Phi P_0\), \(z = \Phi^Tx\), and \(P = P_\epsilon\) with \(\Omega_r^2 = S\).
\end{lemma}
\begin{proof}
\(P_0 = \mathrm{blkdiag}(S, I)\) and \(P_0^{-1}\Phi^T\Pf = \mathrm{blkdiag}(S^{-1}\Phi_q^T\Omega_{\mathrm{full}}^2, \Phi_q^T)\), which gives \(P_\perp\). Multiplying \(\Phi^T\Pfe = \begin{bsmallmatrix}\Phi_q^T\Omega_{\mathrm{full}}^2 & \epsilon\Phi_q^T \\ \epsilon\Phi_q^T & \Phi_q^T\end{bsmallmatrix}\) by \(P_\perp\): the \((1,1)\) block is \(\Phi_q^T\Omega_{\mathrm{full}}^2 - SS^{-1}\Phi_q^T\Omega_{\mathrm{full}}^2 = 0\), the \((1,2)\) and \((2,2)\) blocks are multiples of \(\Phi_q^T(I - \Phi_q\Phi_q^T) = 0\), and the \((2,1)\) block is \(\epsilon\Phi_q^T(I - \Phi_qS^{-1}\Phi_q^T\Omega_{\mathrm{full}}^2) = \epsilon(\Phi_q^T - S^{-1}\Phi_q^T\Omega_{\mathrm{full}}^2)\). It vanishes iff \(\Phi_q^T\Omega_{\mathrm{full}}^2 = S\Phi_q^T\), i.e.\ (transposing) \(\Omega_{\mathrm{full}}^2\Phi_q = \Phi_qS\). Conversely, if \(\Omega_{\mathrm{full}}^2\Phi_q = \Phi_qM\) for some \(M\), then \(S = \Phi_q^T\Phi_qM = M\). The final statements follow by substitution; diagonalizing \(S\) by an orthogonal change of reduced coordinates gives \(\Omega_r\) diagonal.
\end{proof}

\emph{Necessity of (H2b).} In class (H), \(\Af = \begin{bsmallmatrix} 0 & I \\ -\Omega_{\mathrm{full}}^2 & -R_p\end{bsmallmatrix}\) and block multiplication gives \(-(\Pfe\Af + \Af^T\Pfe) = \begin{bsmallmatrix} 2\epsilon\Omega_{\mathrm{full}}^2 & \epsilon R_p \\ \epsilon R_p & 2(R_p - \epsilon I)\end{bsmallmatrix}\). Evaluating \eqref{eq:tail_dissipation} at \(x_\perp = (0, p_\perp)\) with \(p_\perp \in \mathrm{Im}(I - \Phi_q\Phi_q^T)\) (an element of \(\mathcal{V}_\perp\) under (H2a)) gives \(2p_\perp^T(R_p - \epsilon I)p_\perp \ge 2\lambda_{K+1}\|p_\perp\|_2^2 > 0\), which is (H2b).

\subsection*{C. Proof of Lemma~\ref{lem:sector} and Proposition~\ref{prop:deadzone_stability} (Sector Routing and Deadzone Practical Stability)}
\textbf{Sector alignment (Lemma~\ref{lem:sector}).} To verify the constructive sector bound for the \(L\)-layer origin-preserving (\(b_\ell = 0\)) Teacher acting on the port vector \(v := g_r^T P z \in \mathbb{R}^m\) (which equals \(z \in \mathbb{R}^K\) whenever \(m = K\) and \(g_r^T P = I_K\)), let \(W_1 = \cos(\theta_{\mathrm{GS}})U_{\mathrm{sum}}W_{\mathrm{in}} + \sin(\theta_{\mathrm{GS}})U_{\mathrm{diff}}Q_{\mathrm{mix}}W_{\mathrm{in}} \in \mathbb{R}^{H \times m}\) with \(W_{\mathrm{in}}^T W_{\mathrm{in}} = I_m\), let intermediate hidden layers \(W_\ell = U_{\mathrm{sum}}Q_{\mathrm{sum},\ell}U_{\mathrm{sum}}^T + U_{\mathrm{diff}}Q_{\mathrm{diff},\ell}U_{\mathrm{diff}}^T \in O(H)\) (\(2 \le \ell \le L-1\), with \(Q_{\mathrm{sum},\ell}, Q_{\mathrm{diff},\ell} \in O(H/2)\)) be block-diagonal orthogonal, and set \(W_L = -L_{\mathcal{T}}(W_{L-1}\cdots W_1)^T \in \mathbb{R}^{m \times H}\). Since \(\max(a,b) + \min(a,b) = a + b\) and \(\max(a,b) - \min(a,b) = |a - b|\), the orthonormal pairwise sum and difference matrices \(U_{\mathrm{sum}}, U_{\mathrm{diff}}\) satisfy \(U_{\mathrm{sum}}^T \sigma_{\mathrm{GS2}}(y) = U_{\mathrm{sum}}^T y\) and \(\|U_{\mathrm{diff}}^T \sigma_{\mathrm{GS2}}(y)\|_2 = \|U_{\mathrm{diff}}^T y\|_2\) for all \(y \in \mathbb{R}^H\). By induction across \(\ell = 1, \dots, L-1\), the final hidden post-activation \(s_{L-1}(v) := \sigma_{\mathrm{GS2}}(y_{L-1}(v))\) satisfies \(U_{\mathrm{sum}}^T s_{L-1}(v) = \cos(\theta_{\mathrm{GS}})Q_{\mathrm{sum},L-1}\cdots Q_{\mathrm{sum},2}W_{\mathrm{in}} v\) and \(\|U_{\mathrm{diff}}^T s_{L-1}(v)\|_2 = \sin(\theta_{\mathrm{GS}})\|v\|_2\). Defining \(\tilde{Q}_{\mathrm{diff}} := Q_{\mathrm{diff},L-1}\cdots Q_{\mathrm{diff},2}Q_{\mathrm{mix}} \in O(H/2)\), decomposing \(W_{L-1}\cdots W_1v\) in the orthonormal basis \([U_{\mathrm{sum}}\ U_{\mathrm{diff}}]\), and applying Cauchy--Schwarz in \(\mathbb{R}^{H/2}\) yields, under (H-sat) (so that \(\Pi_U\) acts as the identity and \(\Tt(z) = \mathcal{N}(v)\)),
\[
\begin{aligned}
-v^T \mathcal{N}(v) &= L_{\mathcal{T}}\Big(\cos^2(\theta_{\mathrm{GS}})\|W_{\mathrm{in}} v\|_2^2 + \sin(\theta_{\mathrm{GS}})(\tilde{Q}_{\mathrm{diff}}W_{\mathrm{in}} v)^T U_{\mathrm{diff}}^T s_{L-1}(v)\Big)\\ &\ge L_{\mathcal{T}}\big(\cos^2\theta_{\mathrm{GS}} - \sin^2\theta_{\mathrm{GS}}\big)\|v\|_2^2 = \eta_{\mathcal{T}}\|g_r^T P z\|_2^2 \ge 0,
\end{aligned}
\]
and \(-z^TPg_r\Tt(z) = -v^T\mathcal{N}(v)\). Without (H-sat), \(\Pi_U(a) = \lambda a\) with \(\lambda = \min(1, u_{\max}/\|a\|_2) \in (0,1]\) and \(\|\mathcal{N}(v)\|_2 \le L_{\mathcal{T}}\|v\|_2\), so \(-v^T\Pi_U(\mathcal{N}(v)) = -\lambda v^T\mathcal{N}(v) \ge \eta_{\mathcal{T}}\min(1, u_{\max}/(L_{\mathcal{T}}\|v\|_2))\|v\|_2^2\). \hfill \(\blacksquare\)

\begin{lemma}[\(P\)-norm sector bound of the reduced residual]
\label{lem:sector_fnl}
If \(F_{\mathrm{nl}}\) is \(C^1\) on a neighbourhood of \(\Phi\Crz\), then \(\|F_{\mathrm{nl},r}(z)\|_P \le \Lsec\|z\|_P\) for all \(z \in \Crz\).
\end{lemma}
\begin{proof}
\(\Crz\) is a convex ellipsoid containing \(0\), so the segment \(\{sz : s \in [0,1]\}\) lies in \(\Crz\), and \(F_{\mathrm{nl},r}(0) = 0\). By the integral mean-value theorem, \(F_{\mathrm{nl},r}(z) = \int_0^1\nabla F_{\mathrm{nl},r}(sz)\,z\,ds\), hence \(P^{1/2}F_{\mathrm{nl},r}(z) = \int_0^1 \big(P^{1/2}\nabla F_{\mathrm{nl},r}(sz)P^{-1/2}\big)P^{1/2}z\,ds\) and \(\|F_{\mathrm{nl},r}(z)\|_P \le \sup_{s}\|P^{1/2}\nabla F_{\mathrm{nl},r}(sz)P^{-1/2}\|_2\,\|z\|_P \le \Lsec\|z\|_P\). (Only the sector ratio \(\sup_{z \in \Crz\setminus\{0\}}\|F_{\mathrm{nl},r}(z)\|_P/\|z\|_P \le \Lsec\) is used below, and it may replace \(\Lsec\).)
\end{proof}

\textbf{Proof of Proposition~\ref{prop:deadzone_stability}.} With \(w \equiv 0\), \(x_\perp \equiv 0\), the closed-loop reduced trajectory satisfies \(\dot{z} = A_{r,0} z + F_{\mathrm{nl},r}(z) + g_r \Tt(z) + g_r[u(t) - \Tt(z(t))]\) wherever \(\dot z\) exists; since \(u\) is piecewise constant and the right-hand side is continuous in \(z\), \(z\) is piecewise \(C^1\) and has a right derivative \(\dot z^+(t)\) everywhere. Let \(\rho(t) := \|z(t)\|_P\).
For \(z(t) \neq 0\), \(\|\cdot\|_P\) is differentiable at \(z(t)\) and the upper right Dini derivative satisfies \(D^+\rho(t) = z^TP\dot z^+/\|z\|_P\). By the definition of \(\mu_{2,P}\), \(z^TPA_{r,0}z \le -\mu_R\|z\|_P^2\); by Cauchy--Schwarz in the \(P\)-inner product and Lemma~\ref{lem:sector_fnl}, \(z^TPF_{\mathrm{nl},r}(z) \le \|z\|_P\|F_{\mathrm{nl},r}(z)\|_P \le \Lsec\|z\|_P^2\); by Lemma~\ref{lem:sector}, \(z^T P g_r \Tt(z) \le 0\); and by Cauchy--Schwarz in the \(P\)-metric (\(\|P^{1/2} g_r a\|_2 \le \sqrt{\lambda_{\max}(P)}\|g_r\|_2\|a\|_2\)) and hypothesis (ii),
\[
D^+ \|z\|_P \le -\big(\mu_R - \Lsec\big)\|z\|_P + \sqrt{\lambda_{\max}(P)}\,\|g_r\|_2\,\overline{\varepsilon}_q .
\]
For \(z(t) = 0\): \(A_{r,0}0 = 0\), \(F_{\mathrm{nl},r}(0) = 0\), \(\Tt(0) = 0\), so \(\dot z^+(t) = g_ru(t) = g_r[u(t) - \Tt(0)]\) and \(D^+\rho(t) = \limsup_{h\downarrow0}\|z(t+h)\|_P/h = \|\dot z^+(t)\|_P \le \sqrt{\lambda_{\max}(P)}\|g_r\|_2\overline{\varepsilon}_q\), which is the same bound at \(\rho = 0\). This proves (a); no case distinction for the deadzone \(\mathcal{B}_0\) is needed because \(\overline{\varepsilon}_q\) already covers \(\Psi_A(0) = 0\) (see \eqref{eq:eps_A_triangle}).

(b) Let \(c := \max(\rho(t_0), r_{\mathrm{ss}})\) and suppose \(\rho(t_1) > c\) for some \(t_1 > t_0\) in \(I\). Let \(t^\sharp := \sup\{t \in [t_0, t_1] : \rho(t) \le c\}\); by continuity \(\rho(t^\sharp) = c\) and \(\rho > c \ge r_{\mathrm{ss}}\) on \((t^\sharp, t_1]\), where (a) gives \(D^+\rho < 0\). A continuous function with \(D^+\rho \le 0\) on an interval is non-increasing there, so \(\rho(t_1) \le \rho(t^\sharp) = c\), a contradiction. Strict decrease when \(\rho > r_{\mathrm{ss}}\) follows from \(D^+\rho < 0\) by the same monotonicity lemma.

(c) Suppose a periodic solution of period \(T_{\mathrm{per}} > 0\) lies in \(\Crz\) and \(\rho(t_0) > r_{\mathrm{ss}}\) for some \(t_0\). If \(\rho > r_{\mathrm{ss}}\) on \([t_0, t_0 + T_{\mathrm{per}}]\), strict decrease gives \(\rho(t_0 + T_{\mathrm{per}}) < \rho(t_0)\); otherwise \(\rho(t') \le r_{\mathrm{ss}}\) for some \(t' \in (t_0, t_0 + T_{\mathrm{per}}]\), and (b) applied from \(t'\) gives \(\rho(t_0 + T_{\mathrm{per}}) \le r_{\mathrm{ss}} < \rho(t_0)\). Both contradict periodicity.

Validity of (ii) under Theorem~\ref{thm:main_cbf}: by the \(e_u\) bounds of Step~2 in Appendix~D and Lipschitz continuity of \(\Tt\),
\[
\|u(t) - \Tt(z_k)\|_2 \le \max(\eA, \eC + \dul), \qquad \|\Tt(z_k) - \Tt(z(t))\|_2 \le \LTz\|z(t) - z_k\|_2 \le \LTz M_zT_s = \dul
\]
for \(t \in [t_k, t_{k+1})\), so \(\|u(t) - \Tt(z(t))\|_2 \le \overline{\varepsilon}_q\) for all \(t \ge 0\).

Finally, under (H-dz), \(\|z\|_P \le r_{\mathrm{ss}} < r_{\min}\) gives \(\tilde{h}_r(z) = r_{\min}^2 + h_{\mathrm{act}} - \|z\|_P^2 > h_{\mathrm{act}}\), and \(h(\Phi z) > \epsP^2\), so \(z \in \Bint\). \hfill \(\blacksquare\)

\subsection*{D. Proof of Theorem~\ref{thm:main_cbf} and Corollary~\ref{cor:stochastic_w2}}
\textbf{Step 0 (Existence and joint first exit).} The closed-loop right-hand side \(F(x) + Gu(t) + Dw(t)\) is locally Lipschitz in \(x\) (Assumption~\ref{assump:generic}), and measurable and locally bounded in \(t\) (\(u\) is piecewise constant and determined causally by the samples \(z(x(t_k))\); \(w\) is measurable and bounded). By the Carath\'eodory existence and uniqueness theorem, applied successively on \([t_k, t_{k+1})\), there is a unique absolutely continuous solution on a maximal interval \([0, t_{\max})\), and \(\|x(t)\| \to \infty\) as \(t \uparrow t_{\max}\) if \(t_{\max} < \infty\). At \(t = 0\), (H8) gives \(z(0) \in \tilde{\mathcal{C}}_r \subseteq \mathcal{C}_r \subset \mathrm{int}(\Crz)\) and, by Lemma~\ref{lem:quad_split}, \(h(x(0)) = h(\Phi z(0)) - \|x_\perp(0)\|_{\Pfe}^2 \ge \tilde h_r(z(0)) + \Delta h_{\mathrm{dip}} \ge 0\), so \(x(0) \in \mathcal{C} \subset \mathrm{int}(\XO)\) (Lemma~\ref{lem:odd_containment}). Define
\[
t^* := \sup\{t \in [0, t_{\max}) : x(s) \in \XO \text{ and } z(s) \in \Crz \text{ for all } s \in [0, t]\} > 0 .
\]
On \([0, t^*)\), every supremum defining \(M_\perp\), \(\Lperp\) (note \(z \in \Crz \Rightarrow \Phi z \in \mathcal{C} \subset \XO\) by Lemma~\ref{lem:odd_containment}), \(M_z\), \(L_h\), \(\overline g_h\), \(\overline d_h\), \(L_{\mathcal{L}_{f_r}\tilde h_r}\), \(L_{\mathcal{L}_{g_r}\tilde h_r}\) applies, because \(u(t) \in U_{\mathrm{adm}}\) (\(\Psi_A, \Psi_B\) take values in \(U_{\mathrm{adm}}\)) and \(w(t) \in \mathcal{B}(\varepsilon_{\mathrm{adv}})\). Steps 1--2 below, carried out on \([0, t^*)\), give \(\|x_\perp(t)\|_{\Pfe} \le \epsP\), \(z(t) \in \mathcal{C}_r\) and \(x(t) \in \mathcal{C}\) for all \(t \in [0,t^*)\). Suppose \(t^* < t_{\max}\). By continuity and closedness, \(x(t^*) \in \mathcal{C} \subset \mathrm{int}(\XO)\) and \(z(t^*) \in \mathcal{C}_r \subset \mathrm{int}(\Crz)\), so there is \(\delta > 0\) with \(x(s) \in \XO\) and \(z(s) \in \Crz\) on \([0, t^* + \delta]\), contradicting the definition of \(t^*\). Hence \(t^* = t_{\max}\), \(x(t)\) stays in the compact set \(\mathcal{C}\) on \([0, t_{\max})\), and therefore \(t_{\max} = \infty\).

\textbf{Step 1 (Uniform Spectral Tail Bound in the Unified \(\Pfe\)-Norm):} Using \(\Pfe\)-orthogonality (\(\Phi^T \Pfe P_\perp = 0 \implies P_\perp^T \Pfe x_\perp = \Pfe x_\perp\), because \(P_\perp^T\Pfe x_\perp = \Pfe x_\perp - \Pf\Phi P_0^{-1}(\Phi^T\Pfe x_\perp)\) and \(\Phi^T\Pfe x_\perp = \Phi^T\Pfe P_\perp x = 0\)) to project \eqref{eq:full_system} (with \(\Sigma \equiv 0\) and \(\Af = (J_{\mathrm{full}} - R_{\mathrm{full}})\Pf\)) onto \(\mathcal{V}_\perp = \mathrm{Im}(P_\perp)\), i.e.\ \(\dot x_\perp = P_\perp\dot x = P_\perp\Af x_\perp + P_\perp[\Af\Phi z + F_{\mathrm{nl}}(x) + Gu + Dw]\), and differentiating the unified tail Lyapunov function \(V_\perp(x_\perp) := \frac{1}{2}\|x_\perp(t)\|_{\Pfe}^2 = \frac{1}{2}x_\perp^T \Pfe x_\perp\) (absolutely continuous in \(t\)) gives, for a.e.\ \(t\),
\[
\frac{d}{dt}V_\perp(x_\perp) = \frac{1}{2}x_\perp^T \Big(\Pfe \Af + \Af^T \Pfe\Big) x_\perp + x_\perp^T \Pfe P_\perp\Big(\Af\Phi z + F_{\mathrm{nl}}(x) + Gu + Dw\Big).
\]
The quadratic term is bounded by \eqref{eq:tail_dissipation}, i.e.\ hypothesis (H2c) (in class (D), skew-symmetry \(J_{\mathrm{full}}^T = -J_{\mathrm{full}}\) yields \((\Pf x_\perp)^T J_{\mathrm{full}} (\Pf x_\perp) = 0\) identically on \(\mathcal{V}_\perp\) without requiring \(J_{\mathrm{full}}\) and \(R_{\mathrm{full}}\) to commute, so (H2c) reduces to a condition on \(R_{\mathrm{full}}\); in class (H) the rate is computed by Lemma~\ref{lem:pencil_rate}). Applying \eqref{eq:tail_dissipation} and Cauchy--Schwarz in the \(\Pfe\)-inner product together with \eqref{eq:M_perp} yields
\begin{multline*}
\|x_\perp\|_{\Pfe}\frac{d}{dt}\|x_\perp\|_{\Pfe} \le -\lambda_{K+1}\|x_\perp\|_{\Pfe}^2 + M_\perp \|x_\perp\|_{\Pfe} \;\implies\; \\ \|x_\perp(t)\|_{\Pfe} \le \max\!\left(\|x_\perp(0)\|_{\Pfe},\, \frac{M_\perp}{\lambda_{K+1}}\right) \quad \forall t \in [0,t^*).
\end{multline*}
(The implication is the comparison lemma applied to \(V_\perp\): with \(c := \max(\|x_\perp(0)\|_{\Pfe}, M_\perp/\lambda_{K+1})\), if \(\|x_\perp(t_1)\|_{\Pfe} > c\), then on the last interval \((t^\sharp, t_1]\) on which \(\|x_\perp\|_{\Pfe} > c\) we have \(\dot V_\perp \le \|x_\perp\|_{\Pfe}(-\lambda_{K+1}\|x_\perp\|_{\Pfe} + M_\perp) < 0\) a.e., hence \(V_\perp(t_1) < V_\perp(t^\sharp) = c^2/2\), a contradiction.) By (H8), \(\|x_\perp(t)\|_{\Pfe} \le \epsP\) on \([0, t^*)\).

\textbf{Step 2 (Exact Exponentially Weighted Integration Over \([t_k, t_k + T_s]\)):} Let \(s := t - t_k \in [0, T_s)\) and suppose by induction that \(z_k \in \tilde{\mathcal{C}}_r \subseteq \mathcal{C}_r = \Cact \cup \Bint\). Along \eqref{eq:reduced_dynamics}, \(\|\Delta_\perp(z(t_k + s), x_\perp(t_k + s))\|_2 \le \Lperp\epsP\) by \eqref{eq:Lperp_def} and Step 1, hence \(\|\dot z\|_2 \le M_z\) and \(\|z(t_k + s) - z_k\|_2 \le M_z s \le M_z T_s\). Since \(\Crz\) is convex and \(\|\nabla\tilde h_r\|_2 = \|2Pz\|_2 \le L_h\) there, \(\tilde h_r\) is \(L_h\)-Lipschitz on \(\Crz\). If \(z_k \in \Bint\) (so \(\tilde{h}_r(z_k) > h_{\mathrm{act}} > L_h M_z T_s\)), then \(L_h\)-Lipschitz continuity of \(\tilde{h}_r\) immediately ensures \(\tilde{h}_r(z(t_k + s)) \ge \tilde{h}_r(z_k) - L_h M_z s > h_{\mathrm{act}} - L_h M_z T_s > 0\) for all \(s \in [0, T_s]\), so \(z(t_{k+1}) \in \tilde{\mathcal{C}}_r\) and \(\tilde{h}_r(z(t_k + s)) > 0 \ge -\Delta h_{\mathrm{dip}}\) hold automatically. Otherwise, \(z_k \in \Cact \cap \tilde{\mathcal{C}}_r\), where the certified Teacher condition \eqref{eq:teacher_cbf} holds at \(z_k\) with \(\delta_{\mathrm{cert}} \ge \delta_{\mathrm{total}}\). By the chain rule (\(\tilde h_r \in C^1\), \(z\) absolutely continuous), for a.e.\ \(s\), \(\dot{\tilde h}_r(z) = \mathcal{L}_{f_r}\tilde h_r(z) + \mathcal{L}_{g_r}\tilde h_r(z)u + \mathcal{L}_{d_r}\tilde h_r(z)w + \nabla\tilde h_r(z)^T\Delta_\perp\), with \(|\mathcal{L}_{d_r}\tilde h_r(z)w| \le \overline d_h\varepsilon_{\mathrm{adv}}\) and \(|\nabla\tilde h_r^T\Delta_\perp| \le L_h\Lperp\epsP\). Adding and subtracting \eqref{eq:teacher_cbf} evaluated at \(z_k \in \Cact\), writing \(\mathcal{L}_{g_r}\tilde h_r(z)u - \mathcal{L}_{g_r}\tilde h_r(z_k)\Tt(z_k) = [\mathcal{L}_{g_r}\tilde h_r(z) - \mathcal{L}_{g_r}\tilde h_r(z_k)]u + \mathcal{L}_{g_r}\tilde h_r(z_k)[u - \Tt(z_k)]\), and using the Lipschitz constants on \(\Crz\), \(\|u\|_2 \le u_{\max}\) and \(\|z - z_k\|_2 \le M_zs\) (so that the three Lipschitz shifts sum to \(L_{\mathrm{hold}}s\)) gives, for a.e.\ \(s \in [0, T_s)\),
\begin{equation}
\label{eq:ode_diff_ineq}
\dot{\tilde{h}}_r(z(t_k + s)) + \gamma_0 \tilde{h}_r(z(t_k + s)) \ge \delta_{\mathrm{total}} - L_h \Lperp\epsP - \overline{d}_h \varepsilon_{\mathrm{adv}} - L_{\mathrm{hold}}\,s - \overline{g}_h \|u(t_k + s) - \Tt(z_k)\|_2.
\end{equation}
Under the multi-rate schedule \eqref{eq:multirate_hold}, the control mismatch \(e_u(t_k + s) := \|u(t_k + s) - \Tt(z_k)\|_2\) on each sub-interval satisfies:
\begin{itemize}
    \item For \(s \in [0, \tau_A)\): at \(k = 0\), the standby initialization \(u_{B,-1} = \Psi_B(z_0)\) gives \(e_u(s) = \|\Psi_B(z_0) - \Tt(z_0)\|_2 \le \varepsilon_{\mathrm{CROWN}}(b_B) \le \varepsilon_{\mathrm{CROWN}}(b_B) + \Delta u_{\mathrm{lag}}\); for \(k \ge 1\), \(u(t_k + s) = \Psi_B(z_{k-1})\), so \(e_u(t_k + s) \le \|\Psi_B(z_{k-1}) - \Tt(z_{k-1})\|_2 + \|\Tt(z_{k-1}) - \Tt(z_k)\|_2 \le \varepsilon_{\mathrm{CROWN}}(b_B) + \Delta u_{\mathrm{lag}}\), where \(\Delta u_{\mathrm{lag}} = \LTz M_z T_s\).
    \item For \(s \in [\tau_A, \tau_B)\): \(u(t_k + s) = \Psi_A(Q_{b_A}(z_k))\), so \(e_u(t_k + s) \le \varepsilon_A(b_A)\) by \eqref{eq:eps_A_triangle} (including the deadzone case \(Q_{b_A}(z_k) = 0\)); under the optional lower bound of \eqref{eq:codesign_window}, moreover \(\varepsilon_A(b_A) < \varepsilon_{\mathrm{CROWN}}(b_B) + \Delta u_{\mathrm{lag}}\).
    \item For \(s \in [\tau_B, T_s)\): \(u(t_k + s) = \Psi_B(z_k)\), so \(e_u(t_k + s) \le \varepsilon_{\mathrm{CROWN}}(b_B)\).
\end{itemize}
Multiplying \eqref{eq:ode_diff_ineq} by the integrating factor \(e^{\gamma_0 s} > 0\) and integrating from \(s = 0\) to \(s = T_s\) (the left side is \(\frac{d}{ds}[e^{\gamma_0s}\tilde h_r(z(t_k+s))]\) a.e., and \(e^{\gamma_0s}\tilde h_r(z(t_k+s))\) is absolutely continuous) yields
\begin{equation}
\label{eq:exact_exp_integral}
e^{\gamma_0 T_s}\tilde{h}_r(z(t_{k+1})) - \tilde{h}_r(z(t_k)) \ge \int_0^{T_s} e^{\gamma_0 s} \Big( \delta_{\mathrm{total}} - L_h \Lperp\epsP - \overline{d}_h \varepsilon_{\mathrm{adv}} - L_{\mathrm{hold}}\,s - \overline{g}_h\,e_u(t_k + s) \Big) ds.
\end{equation}
Dividing \eqref{eq:exact_exp_integral} by \(\int_0^{T_s} e^{\gamma_0 s}\,ds = \frac{e^{\gamma_0 T_s} - 1}{\gamma_0} > 0\), using the exact sub-interval weights \(w_1(\gamma_0), w_2(\gamma_0), w_3(\gamma_0)\) from \eqref{eq:exp_weights} with \(w_1(\gamma_0) + w_2(\gamma_0) + w_3(\gamma_0) = 1\), and evaluating the exact integration-by-parts identity
\[
\frac{\gamma_0}{e^{\gamma_0 T_s} - 1}\int_0^{T_s} s\,e^{\gamma_0 s}\,ds = \frac{(\gamma_0 T_s - 1)e^{\gamma_0 T_s} + 1}{\gamma_0(e^{\gamma_0 T_s} - 1)} = \varphi(\gamma_0 T_s)\,T_s,
\]
we obtain
\[
\frac{\gamma_0}{e^{\gamma_0 T_s} - 1}\int_0^{T_s} e^{\gamma_0 s}\,e_u(t_k + s)\,ds \le w_2(\gamma_0)\,\varepsilon_A(b_A) + \big(1 - w_2(\gamma_0)\big)\varepsilon_{\mathrm{CROWN}}(b_B) + w_1(\gamma_0)\,\Delta u_{\mathrm{lag}}.
\]
Substituting \(\delta_{\mathrm{total}}\) from \eqref{eq:delta_total_exact} into \eqref{eq:exact_exp_integral} makes the right-hand side of \eqref{eq:exact_exp_integral} non-negative, proving \(e^{\gamma_0 T_s}\tilde{h}_r(z(t_{k+1})) \ge \tilde{h}_r(z(t_k))\) and hence, by induction on \(k \in \mathbb{N}_0\), \(\tilde{h}_r(z(t_k)) \ge 0\) (\(z(t_k) \in \tilde{\mathcal{C}}_r\)) for all \(k \in \mathbb{N}_0\) with \(t_k < t^*\).

For any intermediate continuous time \(t = t_k + \tau\) with \(\tau \in (0, T_s)\), multiplying \eqref{eq:ode_diff_ineq} by \(e^{\gamma_0 s}\), integrating from \(s = 0\) to \(s = \tau\), and dividing by \(e^{\gamma_0 \tau}\) gives
\[
\tilde{h}_r(z(t_k + \tau)) \ge e^{-\gamma_0 \tau}\tilde{h}_r(z(t_k)) + \int_0^\tau e^{-\gamma_0(\tau - s)}\,\mathrm{RHS}(s)\,ds \ge -\int_0^\tau \big[\mathrm{RHS}(s)\big]_-\,ds,
\]
where \([a]_- := \max(0, -a) \ge 0\) and we used \(\tilde{h}_r(z(t_k)) \ge 0\) and \(0 < e^{-\gamma_0(\tau - s)} \le 1\). Because \(\delta_{\mathrm{total}} \ge L_h \Lperp\epsP + \overline{d}_h \varepsilon_{\mathrm{adv}} + \overline{g}_h \varepsilon_{\mathrm{CROWN}}(b_B)\) (as \(\eA \ge \eC\)), the instantaneous deficit \([\mathrm{RHS}(s)]_-\) is bounded by \(\overline{g}_h \Delta u_{\mathrm{lag}} + L_{\mathrm{hold}} s\) on \([0, \tau_A)\), by \(\overline{g}_h(\varepsilon_A(b_A) - \varepsilon_{\mathrm{CROWN}}(b_B)) + L_{\mathrm{hold}} s\) on \([\tau_A, \tau_B)\), and by \(L_{\mathrm{hold}} s\) on \([\tau_B, T_s)\). Integrating from \(s = 0\) to \(\tau < T_s\) yields \(\int_0^\tau [\mathrm{RHS}(s)]_-\,ds \le \Delta h_{\mathrm{dip}}\). Consequently, \(\tilde{h}_r(z(t)) \ge -\Delta h_{\mathrm{dip}}\) (\(z(t) \in \mathcal{C}_r\)) holds for all \(t \in [0, t^*)\), and by the quadratic splitting identity (Lemma~\ref{lem:quad_split}) and Step 1,
\[
h(x(t)) = h\big(\Phi z(t) + x_\perp(t)\big) = h(\Phi z(t)) - \|x_\perp(t)\|_{\Pfe}^2 \ge h(\Phi z(t)) - \epsP^2 = \tilde{h}_r(z(t)) + \Delta h_{\mathrm{dip}} \ge 0 \quad \forall t \in [0, t^*),
\]
ensuring \(x(t) \in \mathcal{C} \subset \XO\). Together with Step 0 (\(t^* = \infty\)) this proves conclusions (a)--(c) of Theorem~\ref{thm:main_cbf}. \hfill \(\blacksquare\)

\textbf{Step 3 (Proof of Corollary~\ref{cor:stochastic_w2}).} Write \(\lambda := \lambda_{K+1}\), \(r_t := \|x_\perp(t)\|_{\Pfe}\), \(\varepsilon_\sigma := \epsS(K,\alpha,T)\), \(c_0 := \epsP(T)\), so \(\varepsilon_\sigma = c_0 + \Delta_\sigma\). All processes live on a filtered probability space satisfying the usual conditions; the strong solution exists up to its explosion time \(\zeta\).

\emph{(a) Disturbance bound.} For any unit \(a \in \mathbb{R}^{n_w}\), \(\xi \mapsto a^T\psi(\xi)\) is \(1\)-Lipschitz and integrable under \(\mu_t, \mu_0 \in \mathcal{P}_2(\Xi)\), so Kantorovich--Rubinstein duality gives \(a^T\int\psi\,d(\mu_t - \mu_0) \le W_1(\mu_t,\mu_0)\); taking the supremum over \(a\) and using \(\int\psi\,d\mu_0 = 0\) and \(W_1 \le W_2\) (Jensen), \(\|\int\psi\,d\mu_t\|_2 \le W_2(\mu_t, \mu_0) \le r_{W_2}t\). Hence \(\|w(t)\|_2 \le \varepsilon_{\mathrm{adv}} + r_{W_2}t \le \varepsilon_{\mathrm{adv}}(T)\) for \(t \le T\), deterministically.

\emph{(b) Stopping times.} Let \(B_t := h(\Phi z_t) - \varepsilon_\sigma^2 = \tilde h_r(z_t) + \Delta h_{\mathrm{dip}}\), \(\tau_{\mathrm{exit}} := \inf\{t \ge 0 : B_t \le 0\}\), \(\tau_\perp := \inf\{t \ge 0 : r_t > \varepsilon_\sigma\}\), \(\tau_{\mathrm{BM}}\) as in (d), and \(\tau_{\mathrm{stop}} := \tau_{\mathrm{exit}} \wedge \tau_{\mathrm{BM}} \wedge \tau_\perp\). For \(t \le \tau_{\mathrm{exit}}\wedge\tau_\perp\) (by continuity of paths, at the stopping time as well), \(B_t \ge 0\) and \(r_t \le \varepsilon_\sigma\), so Lemma~\ref{lem:quad_split} gives \(h(x_t) = h(\Phi z_t) - r_t^2 \ge B_t \ge 0\): \(x_t \in \mathcal{C}\) and \(z_t \in \Crz\). In particular \(\tau_{\mathrm{exit}}\wedge\tau_\perp < \zeta\) (the compact set \(\mathcal{C}\) must be left before explosion), every constant of Section~\ref{sec:theory} (with the replacements of the corollary) applies before \(\tau_{\mathrm{stop}}\), and \(\{\inf_{[0,T]}h(x_t) < 0\} \subseteq \{\tau_{\mathrm{exit}}\wedge\tau_\perp \le T\}\).

\emph{(c) Tail envelope via a generalized It\^o formula.} The tail obeys \(dx_\perp = P_\perp[\Af x + F_{\mathrm{nl}}(x) + Gu + Dw]dt + \Sigma_\perp(x)\,dW_t\). For the \(C^2\) function \(V := \frac12 r^2 = \frac12x_\perp^T\Pfe x_\perp\), It\^o's formula, Step 1, the bound \(\|P_\perp D(w - w_{\mathrm{adv}})\|_{\Pfe} \le \|\Pfe^{1/2}P_\perp D\|_2r_{W_2}T\) from (a), and \(\frac12\mathrm{tr}(\Sigma_\perp^T\Pfe\Sigma_\perp) = \frac12\|\Pfe^{1/2}P_\perp\Sigma\|_F^2 \le \frac12\overline\sigma_\perp^2\) give, for \(t < \tau_{\mathrm{stop}} \wedge T\),
\(dV \le (-2\lambda V + M_\perp(T)\sqrt{2V} + \tfrac12\overline\sigma_\perp^2)dt + dN_t\), \(dN_t := x_\perp^T\Pfe\Sigma_\perp\,dW_t\), with \(d\langle N\rangle_t = \|\Sigma_\perp^T\Pfe x_\perp\|_2^2dt \le \|\Pfe^{1/2}P_\perp\Sigma\|_2^2\,r_t^2dt \le \overline\sigma_\perp^2r_t^2dt\) (\(\|\cdot\|_2 \le \|\cdot\|_F\)). Let \(g(V) := (\sqrt{2V} - c_0)_+^2\), so \(U_t := g(V_t) = (r_t - c_0)_+^2\). For \(\sqrt{2V} > c_0\), \(g'(V) = 2(1 - c_0/r) \in [0,2)\) and \(g''(V) = 2c_0/r^3 > 0\); \(g' = g'' = 0\) for \(\sqrt{2V} < c_0\); \(g'\) is continuous at \(\sqrt{2V} = c_0\). Thus \(g\) is convex and \(C^{1,1}\) on \([0,\infty)\) (its second derivative is bounded by \(2/c_0^2\) where non-zero), and the Meyer--It\^o formula (Revuz--Yor, Thm.~VI.1.5), whose second-derivative measure \(g''(V)dV\) has no atoms so that the local-time term reduces via the occupation-times formula, gives \(dU = g'(V)dV + \frac12g''(V)d\langle V\rangle\). On \(\{r > c_0\}\), using \(g', g'' \ge 0\), \(M_\perp(T) = \lambda c_0\) and \(d\langle V\rangle = d\langle N\rangle \le \overline\sigma_\perp^2r^2dt\):
\[
\text{drift}(U) \le \tfrac{2(r - c_0)}{r}\big(-\lambda r^2 + \lambda c_0 r + \tfrac12\overline\sigma_\perp^2\big) + \tfrac{c_0}{r^3}\overline\sigma_\perp^2r^2 = -2\lambda(r - c_0)^2 + \overline\sigma_\perp^2 = -2\lambda U + \overline\sigma_\perp^2,
\]
and on \(\{r \le c_0\}\) the drift is \(0 \le -2\lambda U + \overline\sigma_\perp^2\). Hence \(dU \le (-2\lambda U + \overline\sigma_\perp^2)dt + d\tilde M_t\) with \(\tilde M := \int g'(V)dN\), \(d\langle\tilde M\rangle \le \frac{4(r-c_0)^2}{r^2}\overline\sigma_\perp^2r^2dt = 4\overline\sigma_\perp^2U\,dt\). This is the tail SDE with Hilbert--Schmidt norm. Applying the \(C^2\) It\^o formula to \(e^{\eta U}\) with \(\eta = \lambda/(2\overline\sigma_\perp^2)\),
\[
\mathcal{L}e^{\eta U} \le \eta e^{\eta U}\big(-2\lambda U + \overline\sigma_\perp^2 + 2\eta\overline\sigma_\perp^2U\big) = \eta e^{\eta U}\big(\overline\sigma_\perp^2 - \lambda U\big) = \frac{\lambda}{2}e^{\eta U}\big(1 - 2\eta U\big) \le \frac{\lambda}{2},
\]
using \(e^x(1 - 2x) \le e^x(1 - x) \le 1\) for \(x = \eta U \ge 0\). Let \(\sigma := \tau_{\mathrm{stop}}\wedge T\). On \([0,\sigma]\), \(U \le \Delta_\sigma^2\) and \(r \le \varepsilon_\sigma\), so the stochastic integral in \(e^{\eta U_{t\wedge\sigma}}\) has bounded integrand and is a true martingale; optional stopping gives \(\mathbb{E}[e^{\eta U_\sigma}] \le e^{\eta U_0} + \frac{\lambda}{2}\mathbb{E}[\sigma] \le 1 + \frac12\lambda T\), since \(U_0 = 0\) by \(r_0 \le M_\perp(T)/\lambda = c_0\). On \(A_\perp := \{\tau_\perp \le T\wedge\tau_{\mathrm{BM}}\wedge\tau_{\mathrm{exit}}\}\), \(\sigma = \tau_\perp\) and \(U_\sigma = \Delta_\sigma^2\); Markov's inequality yields \(\mathbb{P}(A_\perp) \le (1 + \frac12\lambda T)e^{-\eta\Delta_\sigma^2} = \alpha/4\) by the definition of \(\Delta_\sigma\) (\(\mathbb{P}(A_\perp) = 0\) when \(\overline\sigma_\perp = 0\), since then \(U \equiv 0\)).

\emph{(d) Brownian excursion tube.} Let \(t_k = \lfloor t/T_s\rfloor T_s\), and define the stopped projected Brownian increment
\[
\xi_k(t) := \int_{t_k}^t \mathbb{I}_{\{\vartheta \le \tau_{\mathrm{exit}}\wedge\tau_\perp\}}\Sigma_r(x_\vartheta)\,dW_\vartheta
\]
and \(\tau_{\mathrm{BM}} := \inf\{t \in [0,T] : \|\xi_k(t)\|_2 > r_{\mathrm{BM}}(\alpha)\}\). Since \(\|\xi_k(t)\|_2 > r_{\mathrm{BM}}(\alpha) \implies \max_{1 \le j \le K} |\xi_{k,j}(t)| > r_{\mathrm{BM}}(\alpha)/\sqrt{K}\) and each coordinate \(\xi_{k,j}(t)\) has quadratic variation \(\langle \xi_{k,j} \rangle_t \le \overline{\sigma}_r^2 T_s\) on \([t_k, t_{k+1}]\) (the integrand is stopped while \(x \in \mathcal{C}\), where \(\|\Sigma_r\|_2 \le \overline\sigma_r\)), the Dambis--Dubins--Schwarz exponential martingale inequality and union bound across \(K N_T\) events yield
\[
\mathbb{P}(\tau_{\mathrm{BM}} \le T) \le 2 K N_T \exp\!\left( -\frac{r_{\mathrm{BM}}(\alpha)^2}{2 K \overline{\sigma}_r^2 T_s} \right) = 2 K N_T \cdot \frac{\alpha}{8 K N_T} = \frac{\alpha}{4}.
\]
For \(t < \tau_{\mathrm{stop}}\), \(z_t - z_{t_k} = \int_{t_k}^t[f_r + g_ru + d_rw + \Delta_\perp]\,d\vartheta + \xi_k(t)\), and the drift has norm \(\le M_z(T)\), so \(\|z_t - z_{t_k}\|_2 \le M_z(T)(t - t_k) + r_{\mathrm{BM}}(\alpha)\) and \(\|z_{t_k} - z_{t_{k-1}}\|_2 \le M_z(T)T_s + r_{\mathrm{BM}}(\alpha)\).

\emph{(e) Exponential barrier process.} Let \(\theta := \thK = \ln(2/\alpha)/h_0\) and \(Z_t := \exp(-\theta B_{t\wedge\tau_{\mathrm{stop}}})\). By the \(C^2\) It\^o formula applied to the smooth \(\tilde h_r\) along \(dz = [\cdot]dt + \Sigma_r dW\), \(dB = (\mu_B + \frac12\mathrm{tr}(\Sigma_r^T\nabla_z^2\tilde h_r\Sigma_r))dt + \nabla_z\tilde h_r^T\Sigma_r\,dW\), and applied to \(b \mapsto e^{-\theta b}\), for \(t < \tau_{\mathrm{stop}}\),
\(\mathcal{L}Z_t = -\theta Z_t\big(\mu_B(t) + \frac{1}{2}\mathrm{tr}(\Sigma_r^T\nabla_z^2\tilde{h}_r\Sigma_r) - \frac{1}{2}\theta\|\nabla_z\tilde{h}_r^T\Sigma_r\|_2^2\big)\). We split on the sample point \(z_{t_k} \equiv z_k\), which is \(\mathcal{F}_{t_k}\)-measurable; since \(B_{t_k} > 0\), \(z_{t_k} \in \mathcal{C}_r = \Cact \cup \Bint\).
\begin{enumerate}
    \item \emph{Case A, \(z_{t_k} \in \Cact\)} (\(\tilde h_r(z_{t_k}) \le h_{\mathrm{act}} = \overline h_r\)). As in Step 2, but pointwise in \(t\) and without the factor \(\gamma_0\tilde h_r\) on the left: \(\mu_B(t) \ge \delta_{\mathrm{cert}} - \gamma_0\tilde h_r(z_{t_k}) - (L_{\mathcal{L}_{f_r}\tilde h_r} + L_{\mathcal{L}_{g_r}\tilde h_r}u_{\max})\|z_t - z_{t_k}\|_2 - \overline g_h\|u(t) - \Tt(z_{t_k})\|_2 - \overline d_h\|w(t)\|_2 - L_h\Lperp\varepsilon_\sigma\). By (a), (d), the \(e_u\) bounds of Step 2 with the lag enlarged to \(\LTz(M_z(T)T_s + r_{\mathrm{BM}}) = \dul + \LTz r_{\mathrm{BM}}\), and \(\tilde h_r(z_{t_k}) \le \overline h_r\),
    \[
    \mu_B(t) \ge \delta_{\mathrm{cert}} - \gamma_0\overline h_r - \big[\delta_{\mathrm{total},\sigma} + \Delta\delta_{\mathrm{pw}}\big] - \overline d_hr_{W_2}t \ge \delta_{\mathrm{It\hat o}} + \tfrac12\theta\sigma_h^2 + \big(\delta_{\mathrm{drift}} - \overline d_hr_{W_2}t\big),
    \]
    where we used \(\delta_{\mathrm{total},\sigma} + \Delta\delta_{\mathrm{pw}} = L_h\Lperp\varepsilon_\sigma + \overline g_h\max(\eA, \eC + \dul) + \overline d_h\varepsilon_{\mathrm{adv}} + L_{\mathrm{hold}}T_s + (L_{\mathcal{L}_{f_r}\tilde h_r} + L_{\mathcal{L}_{g_r}\tilde h_r}u_{\max} + \gamma_0L_h + \overline g_h\LTz)r_{\mathrm{BM}}\), which dominates the subtracted terms (the \(\gamma_0L_h\) contributions are slack), and the hypothesis \(\delta_{\mathrm{total},\sigma} + \delta_{\mathrm{stoch}} + \delta_{\mathrm{drift}} \le \delta_{\mathrm{cert}}\) with \eqref{eq:delta_stoch}. Since \(\frac12\mathrm{tr}(\cdot) \ge -\delta_{\mathrm{It\hat o}}\), \(\|\nabla\tilde h_r^T\Sigma_r\|_2 \le \sigma_h\) and \(t \le T \le \delta_{\mathrm{drift}}/(\overline d_hr_{W_2})\), the bracket in \(\mathcal{L}Z_t\) is \(\ge 0\), so \(\mathcal{L}Z_t \le 0\).
    \item \emph{Case B, \(z_{t_k} \in \Bint\)} (\(\tilde h_r(z_{t_k}) > h_{\mathrm{act}}\)). By (d) and \(L_h\)-Lipschitz continuity on \(\Crz\), \(\tilde h_r(z_t) > h_{\mathrm{act}} - L_h(M_z(T)T_s + r_{\mathrm{BM}}) = h_{\mathrm{act}}^-\), so \(B_t > h_{\mathrm{act}}^- + \Delta h_{\mathrm{dip}}\) and \(Z_t < e^{-\theta(h^-_{\mathrm{act}} + \Delta h_{\mathrm{dip}})}\); and \(\mu_B(t) \ge -\underline\mu^-_{\mathrm{int}}\) because \((z_t, u(t), w(t), x_{\perp,t})\) is admissible in the supremum defining \(\underline\mu^-_{\mathrm{int}}\). Hence \(\mathcal{L}Z_t \le \theta Z_t(\underline\mu^-_{\mathrm{int}} + \delta_{\mathrm{It\hat o}} + \frac12\theta\sigma_h^2) \le c_{\mathrm{int}}\).
\end{enumerate}
\emph{(f) Kushner step.} \(Z_{t\wedge\tau_{\mathrm{stop}}} \in (0,1]\) (as \(B \ge 0\) up to \(\tau_{\mathrm{exit}}\)), and its stochastic integral has integrand bounded by \(\theta\sigma_h\), hence is a true martingale. By (e), \(Y_t := Z_{t\wedge\tau_{\mathrm{stop}}} + c_{\mathrm{int}}(T - t\wedge\tau_{\mathrm{stop}})\), \(t \in [0,T]\), is a non-negative continuous supermartingale, and Ville's (Doob's) maximal inequality gives \(\mathbb{P}(\sup_{[0,T]}Y_t \ge 1) \le Y_0 = e^{-\theta\beta_0} + c_{\mathrm{int}}T\). On \(A_{\mathrm{exit}} := \{\tau_{\mathrm{exit}} \le T\wedge\tau_{\mathrm{BM}}\wedge\tau_\perp\}\), \(B_{\tau_{\mathrm{exit}}} = 0\), so \(Y \ge Z = 1\) at \(\tau_{\mathrm{exit}}\). The second bound in \eqref{eq:T_TTL} gives \(1 - 2c_{\mathrm{int}}T/\alpha \ge e^{-\theta(\beta_0 - h_0)} > 0\), i.e.\ \(e^{-\theta\beta_0} \le e^{-\theta h_0}(1 - 2c_{\mathrm{int}}T/\alpha) = \frac{\alpha}{2} - c_{\mathrm{int}}T\) (using \(e^{-\theta h_0} = \alpha/2\)); hence \(\mathbb{P}(A_{\mathrm{exit}}) \le \alpha/2\).

\emph{(g) Union.} If \(\tau_{\mathrm{exit}}\wedge\tau_\perp \le T\), then either \(\tau_{\mathrm{BM}} < \tau_{\mathrm{exit}}\wedge\tau_\perp\), so \(\tau_{\mathrm{BM}} \le T\) (probability \(\le \alpha/4\) by (d)), or \(\tau_{\mathrm{stop}} = \tau_{\mathrm{exit}}\wedge\tau_\perp \le T\) and the event \(A_{\mathrm{exit}}\) or \(A_\perp\) occurs. Therefore, by (b),
\[
\mathbb{P}\Big(\inf_{t \in [0,T]} h(x_t) < 0\Big) \le \mathbb{P}(\tau_{\mathrm{BM}} \le T) + \mathbb{P}(A_{\mathrm{exit}}) + \mathbb{P}(A_\perp) \le \frac{\alpha}{4} + \frac{\alpha}{2} + \frac{\alpha}{4} = \alpha .
\]
Thus \(\mathbb{P}\big(\inf_{t \in [0, T]} h(x(t)) \ge 0\big) \ge 1 - \alpha\). \hfill \(\blacksquare\)
